\begin{lemma}\label{lem:decic242-last-support}
The remaining support in case $26$ cannot be the reduction of a
rational effective divisor of degree $10$ with entirely
noncuspidal support.
\end{lemma}
\begin{proof}
Write its noncuspidal part as $\overline N$.
For an exact identification, let
\[
 k=\mathbb F_5[\alpha]/
 (\alpha^6+\alpha^4+4\alpha^3+\alpha^2+2).
\]
The coefficient vectors of $B,C\in k$, in ascending powers of
$\alpha$, are
\[
 (B_i)=(0,0,2,4,2,3),\qquad(C_i)=(4,1,4,0,0,1).
\]
The level structure
\[
 E:\ y^2=x^3+Cx^2+B(1-B-C)x,\qquad P=(0,0),\quad Q=(B,B)
\]
represents $\overline N$.
Put
\[
 h(P,Q)=(P+21Q,P+4Q).
\]
This automorphism has order $3$ and agrees with the second power
of Frobenius on $\overline N$.
It acts freely on all $72$ geometric cusps.

Let $\pi:X\to W=X/\langle h\rangle$.
The quotient has genus $17$, and
$\overline N=\pi^*q$ for a closed quadratic place $q$.
We describe the model and the section calculation used below.
With the notation of Lemma~\ref{lem:decic242-quotient5}, adjoining
\[
 w^3-\frac{N(t)}{48}w+
 \frac{M(t)(2t+13)}{1728D(t)}-\frac{u}{8D(t)}=0
\]
to the genus-$5$ field gives the field of $W$.
Indeed, if $r_0,r_1,r_2$ are the ordered normalized $2$-division
roots and $v_0=v$, $v_1=(v-1)/v$, $v_2=1/(1-v)$, then
$w=r_0v_0+r_1v_1+r_2v_2$ satisfies this equation.
The generator used for this invariant-field model is
$(21Q,P+37Q)=h^2$; it agrees with the fourth power of Frobenius
on $\overline N$. Thus it gives the same quotient.

Over $\mathbb F_5$, the genus-zero $(u,w)$-field has a parameter
$z$ with
\[
 u=\frac{z^3+3z^2+4z+3}{z^2+1}.
\]
Adjoining the coordinates of the nonzero $3$-point on $E_u$
gives an absolute degree-$8$ model over $\mathbb F_5(z)$ of
genus $17$.
The source function $u$ is regular on $\overline N$ and satisfies
$u^2+4u+2=0$ there.
We retain \emph{all} quadratic places of $W$ above this polynomial.
There are two. Both are unramified with inert cubic fibre, so
they form a containing set for the actual place $q$.
This avoids evaluating a chart with poles at the CM residue point.

The source cusp signature is $1^{18},2^9,3^6,6^3$.
By Lemma~\ref{lem:decic242-case26-cubic}, the complement
$\overline C$ consists of one cubic and one rational cusp.
The quotient cusp signature is $1^6,2^3,3^2,6$.
The $108$ possible source complements give exactly $12$ divisors
$\pi_*\overline C$, one for each choice of a cubic and a rational
cusp of $W$; each choice has nine source preimages.

We compute the required section dimensions after extending
constants to $\mathbb F_{25}$.
Put $s=u^2+8$, $D=s^2-3s+9$, and
$A=2s^3-9s^2+27s-27$.
Adjoining the Tate parameter $v$ recovers the cyclic degree-$3$
source cover:
\[
 v^3-sv^2+(s-3)v+1=0.
\]
For a primitive cube root $\zeta\in\mathbb F_{25}$, this is the
Kummer cover $y^3=f$, where
\[
 f=\frac{A+3(\zeta-\zeta^2)D}{2}.
\]
The identities
$A^2-4D^3=-27D^2$ and $3v-s=y+D/y$ verify the correspondence.
For $j=0,1,2$, put
\[
 D_j=\left\lfloor\frac{j\operatorname{div}(f)}3\right\rfloor.
\]
Their degrees are $0,-6,-6$, and the branch divisor has degree
$12$. The usual integral local bases of a tame Kummer cover give
\[
 h^0_X(\pi^*T)=\sum_{j=0}^2h^0_W(T+D_j).
\]

A quadratic place over $\mathbb F_5$ becomes a sum of two rational
places over $\mathbb F_{25}$.
The degree-$6$ source place becomes a sum of two cubic places;
we retain both summands in each divisor.
For each of the two admitted $q$ and all $12$ cusp choices, take
$T=q+\pi_*\overline C$, of degree $6$.
The $24$ exact calculations, comprising $72$ Riemann--Roch spaces,
give
\[
 (h^0(T+D_0),h^0(T+D_1),h^0(T+D_2))=(1,0,0).
\]
Consequently
\begin{equation}\label{eq:decic242-final-section}
 h^0_X(\overline N+\pi^*\pi_*\overline C)=1
\end{equation}
for every compatible complement.
All places are exported by two exact ideal generators.
The field is an absolute extension of degree $8$ over the rational
function field; its finite and infinite ideal module bases each
have eight elements. A second native calculation reconstructs
the character divisors, compares both module lattices in both
directions, and reproduces all $72$ dimensions.
Independent arithmetic also verifies the cyclic-cover identity,
the Frobenius phases and the complete cusp-orbit coverage.

Suppose that a rational noncuspidal divisor $D$ of degree $10$
has the proposed reduction.
Lift $\overline C$ to a cuspidal divisor $C$ over an unramified
local field at $5$.
The class $[hD-D-hC+C]$ is torsion by rank zero and
Manin--Drinfeld, and reduces to zero.
Odd unramified torsion injectivity makes it zero.
Thus $L=[D-C]$ is $h$-invariant.
Put
\[
 \Omega=C+hC+h^2C=\pi^*\pi_*C.
\]
The effective divisors $D+\Omega-C$ and $hD+\Omega-hC$ are
linearly equivalent.
By \eqref{eq:decic242-final-section} and upper semicontinuity,
their common line bundle has at most one section, so the divisors
are equal. Comparing their cuspidal parts gives $C=hC$.
Since $h$ acts freely on the geometric cusps, this forces
$3\mid\deg C=4$, a contradiction.
\end{proof}
